% Part: incompleteness
% Chapter: introduction
% Section: decidability

\documentclass[../../../include/open-logic-section]{subfiles}

\begin{document}

\olfileid[pt]{inc}{int}{dec}

\olsection{Indecidibilidade e incompletude}

A prova dos teoremas da incompletude de G\"odel exige a aritmetização
da sintaxe. Mas, mesmo sem ela, podemos obter alguns resultados
interessantes apenas com a hipótese de que uma teoria !!{represents}
todas as relações !!{decidable}. A prova é um argumento diagonal
semelhante à prova da indecidibilidade do problema da parada.

\begin{thm}
Se $\Gamma$ é uma teoria consistente que !!{represents} toda relação
!!{decidable}, então $\Gamma$ não é !!{decidable}.
\end{thm}

\begin{proof}
Suponha que $\Gamma$ fosse !!{decidable}. Mostraremos que, se $\Gamma$
!!{represents} toda relação !!{decidable}, então ela deve ser
inconsistente.

Propriedades !!^{decidable} (relações de um lugar) são representadas
por !!{formula}s com uma variável livre. Seja $!A_0(x)$, $!A_1(x)$,
\dots, uma enumeração computável de todas essas !!{formula}s. Considere
agora o seguinte conjunto $D \subseteq \Nat$:
\[
D = \Setabs{n}{\Gamma \Proves \lnot !A_n(\num{n})}
\]
O conjunto $D$ é !!{decidable}, pois podemos testar se $n \in D$
computando primeiro $!A_n(x)$ e, a partir dela,
$\lnot !A_n(\num{n})$. Evidentemente, substituir o termo $\num{n}$ por
toda ocorrência livre de $x$ em $!A_n(x)$ e prefixar
$!A(\num{n})$ por $\lnot$ é uma operação mecânica. Por hipótese,
$\Gamma$ é !!{decidable}, de modo que podemos testar se
$\lnot !A(\num{n}) \in \Gamma$. Se pertencer, $n \in D$; se não
pertencer, $n \notin D$. Logo, $D$ também é !!{decidable}.

Como $\Gamma$ !!{represents} todas as propriedades !!{decidable}, ela
!!{represents}~$D$. E as !!{formula}s que !!{represents} $D$
em~$\Gamma$ estão todas entre $!A_0(x)$, $!A_1(x)$, \dots. Seja, então,
$d$ um número tal que $!A_d(x)$ !!{represents} $D$ em~$\Gamma$. Se
$d \notin D$, então, como $!A_d(x)$ !!{represents}~$D$,
$\Gamma \Proves \lnot !A_d(\num{d})$. Mas isso significa que $d$
satisfaz a condição definidora de~$D$ e, portanto, $d \in D$. Isso
contradiz $d \notin D$. Logo, por prova indireta, $d \in D$.

Como $d \in D$, pela definição de~$D$, $\Gamma \Proves \lnot
!A_d(\num{d})$. Por outro lado, como $!A_d(x)$ !!{represents}~$D$
em $\Gamma$, $\Gamma \Proves !A_d(\num{d})$. Portanto, $\Gamma$ é
inconsistente.
\end{proof}

\begin{explain}
O teorema anterior mostra que nenhuma teoria consistente que
!!{represents} todas as relações !!{decidable} pode ser !!{decidable}.
Mostraremos que $\Th{Q}$ de fato !!{represents} todas as relações
!!{decidable}; isso significa que todas as teorias que incluem
$\Th{Q}$, como $\Th{PA}$ e $\Th{TA}$, também as representam e, por
conseguinte, tampouco são !!{decidable}. (Como todas essas teorias são
verdadeiras no modelo padrão, todas são consistentes.)

Também podemos usar esse resultado para obter uma versão fraca do
primeiro teorema da incompletude. Toda teoria !!{axiomatizable} e
!!{complete} é !!{decidable}. Assim, teorias consistentes que são
!!{axiomatizable} e !!{represents} todas as propriedades !!{decidable}
não podem ser !!{complete}.
\end{explain}

\begin{thm}
Se $\Gamma$ é !!{axiomatizable} e !!{complete}, então é !!{decidable}.
\end{thm}

\begin{proof}
Toda teoria inconsistente é !!{decidable}, pois teorias inconsistentes
contêm todas as !!{sentence}s; assim, a resposta à pergunta ``$!A \in
\Gamma$?'' é sempre ``sim'', O.s.\@, pode ser decidida.

Suponha, então, que $\Gamma$ seja consistente e, além disso,
!!{axiomatizable} e !!{complete}. Como $\Gamma$ é !!{axiomatizable},
ela é !!{computably enumerable}. Com efeito, podemos enumerar por uma
função computável todas as !!{derivation}s corretas a partir dos
axiomas de~$\Gamma$. De uma !!{derivation} correta podemos computar a
!!{sentence} que ela !!{derive}; conjuntamente, portanto, há uma função
computável que enumera todos os teoremas de~$\Gamma$. Uma !!{sentence}
é teorema de~$\Gamma$ sse $\lnot !A$ não é teorema, pois $\Gamma$
é consistente e !!{complete}. Podemos, então, decidir se
$!A \in \Gamma$ da seguinte maneira. Enumeramos todos os teoremas de
$\Gamma$. Quando $!A$ aparece nessa lista, sabemos que
$\Gamma \Proves !A$. Quando $\lnot !A$ aparece nela, sabemos que
$\Gamma \Proves/ !A$. Como $\Gamma$ é !!{complete}, um desses casos
acabará ocorrendo e o procedimento acabará produzindo uma resposta.
\end{proof}

\begin{cor}
\ollabel{cor:incompleteness}
Se $\Gamma$ é consistente, !!{axiomatizable} e !!{represents} toda
propriedade !!{decidable}, então não é !!{complete}.
\end{cor}

\begin{proof}
Se $\Gamma$ fosse !!{complete}, seria !!{decidable} pelo teorema
anterior (pois é !!{axiomatizable} e consistente). Mas, como $\Gamma$
!!{represents} toda propriedade !!{decidable}, ela não é !!{decidable},
pelo primeiro teorema.
\end{proof}

\begin{prob}
Mostre que $\Th{TA} = \Setabs{!A}{\Sat{N}{!A}}$ não é
!!{axiomatizable}. Você pode supor que $\Th{TA}$ representa todas as
propriedades decidíveis.
\end{prob}

Uma vez estabelecido que, p.ex., $\Th{Q}$ !!{represents} todas
as propriedades !!{decidable}, o corolário nos diz que $\Th{Q}$ deve
ser incompleta. Contudo, sua prova não fornece um exemplo de
!!{sentence} independente; apenas mostra que tal !!a{sentence} deve
existir. Para obtê-lo, precisamos aritmetizar a sintaxe e seguir a ideia
original da prova de G\"odel. E, naturalmente, ainda precisamos mostrar
a primeira afirmação: que $\Th{Q}$ de fato !!{represents} todas as
propriedades !!{decidable}.

Convém observar que nem toda teoria \emph{interessante} é incompleta ou
indecidível. Há muitas teorias suficientemente fortes para descrever
fatos matemáticos interessantes que não satisfazem as condições do
resultado de G\"odel. Por exemplo,
$\Th{Pres} = \Setabs{!A \in \Lang{L_{A^+}}}{\Sat{N}{!A}}$, o conjunto
das !!{sentence}s da linguagem da aritmética sem~$\times$ que são
verdadeiras no modelo padrão, é completo e decidível. Essa teoria é
chamada aritmética de Presburger e prova todas as verdades sobre os
números naturais que podem ser formuladas apenas com $\Obj{0}$,
$\prime$ e~$+$.

\end{document}
